% ============================================================
%  Methodology, Results, and Observations
%  ML4NLU Paper — Tigrigna Gender Bias in XLM-RoBERTa
% ============================================================

% ============================================================
\section{Methodology}
% ============================================================

\subsection{Language and Corpus Construction}

We focus on Tigrigna (\texttt{ti}), a Semitic language spoken primarily in
Eritrea and the Tigray region of Ethiopia, written in Ge'ez (Ethiopic) script.
Tigrigna is a genuinely low-resource language: no substantial pretrained
monolingual model exists, and its representation in multilingual corpora is
sparse.

Because no sufficiently large, profession-annotated Tigrigna corpus is publicly
available, we constructed a synthetic corpus of \textbf{14{,}960 sentences}.
The corpus was built using a set of controlled sentence templates, each
instantiating one of 30 target profession words in a short declarative context.
Templates were designed to co-occur the profession word with gender-marked
anchor terms (pronouns, kinship terms, and generic gender nouns) drawn from
Tigrigna morphology. Every profession word appears in exactly 499 sentences,
distributed across male-gendered, female-gendered, and gender-neutral contexts.
Each sentence was automatically labelled as \texttt{M}, \texttt{F}, or
\texttt{N} by checking for the presence of known gender marker tokens
(see Table~\ref{tab:anchors}).

\begin{table}[h]
\centering
\small
\begin{tabular}{llll}
\toprule
\textbf{Gender} & \textbf{Type} & \textbf{Tigrigna} & \textbf{English gloss} \\
\midrule
Male   & Pronoun   & ንሱ, ንዕኡ, ናቱ              & he, him, his \\
Male   & Kinship   & ኣቦ, ሓው, ወዲ ሓው, በዓል ገዛ   & father, brother, nephew, husband \\
Male   & Generic   & ሰብኣይ, ወዲ, ተባዕታይ, ኣቦሓጎ  & man, boy/son, male, grandfather \\
\midrule
Female & Pronoun   & ንሳ, ንዓኣ, ናታ             & she, her, her (poss.) \\
Female & Kinship   & ኣደ, ሓፍቲ, ጓል ሓፍቲ, ሰበይቲ  & mother, sister, niece, wife \\
Female & Generic   & ጓል, ኣንስተይቲ, ዓባየይ        & girl/daughter, female, grandmother \\
\bottomrule
\end{tabular}
\caption{Gender anchor words used to build the gender subspace and classify
corpus sentences. 11 male and 10 female anchors are used.}
\label{tab:anchors}
\end{table}

\subsection{Model}

All experiments use \textbf{XLM-RoBERTa-base}~\cite{conneau2020unsupervised}
(\texttt{xlm-roberta-base}), a 12-layer transformer with 768-dimensional hidden
states trained on 100 languages via masked language modelling on CommonCrawl
data. The model is used strictly in inference mode with no fine-tuning; all
weights are frozen. We extract hidden states from all 13 layers (the embedding
layer~0 plus the 12 transformer blocks), giving a complete picture of how
representations evolve through the network. All experiments run on CPU with
full determinism; no GPU is required.

For any word $w$ that tokenises into $k$ subword pieces
$t_1, \ldots, t_k$ (using the SentencePiece vocabulary of XLM-RoBERTa), we
obtain its representation at layer $l$ as the mean of the $k$ hidden-state
vectors at the corresponding token positions:
\begin{equation}
  \mathbf{s}_l(w) = \frac{1}{k} \sum_{i=1}^{k} \mathbf{h}_l(t_i)
  \label{eq:meanpool}
\end{equation}
When a word appears in multiple sentences, we further average the per-sentence
vectors to produce a single representation per layer.

\subsection{Gender Geometry Construction}

For each transformer layer $l$, we build a \emph{gender subspace} from the
anchor words as follows.

\paragraph{Anchor embeddings.}
For each male anchor word $a^M_j$ (resp.\ female anchor $a^F_j$), we collect
12 corpus sentences containing that word and compute the layer-$l$
representation $\mathbf{s}_l(a^M_j)$ via Equation~\ref{eq:meanpool}.

\paragraph{Gender centroids.}
The male centroid and female centroid at layer $l$ are:
\begin{equation}
  \mathbf{m}_l = \frac{1}{|A^M|} \sum_{j} \mathbf{s}_l(a^M_j), \quad
  \mathbf{f}_l = \frac{1}{|A^F|} \sum_{j} \mathbf{s}_l(a^F_j)
\end{equation}
where $A^M$ and $A^F$ are the sets of male and female anchor words
($|A^M|=11$, $|A^F|=10$).

\paragraph{Gender direction.}
The unit gender direction vector at layer $l$ is:
\begin{equation}
  \mathbf{g}_l = \frac{\mathbf{m}_l - \mathbf{f}_l}
                      {\|\mathbf{m}_l - \mathbf{f}_l\|}
\end{equation}
By construction $\|\mathbf{g}_l\| = 1$; positive projection onto $\mathbf{g}_l$
indicates proximity to the male cluster and negative projection indicates
proximity to the female cluster.

\subsection{Bias Estimation}

For each profession word $p$, we collect 20 corpus sentences and compute the
layer-$l$ representation $\mathbf{s}_l(p)$ via Equation~\ref{eq:meanpool}.
We employ two independent estimators of gender bias:

\paragraph{Projection bias (ProjBias).}
\begin{equation}
  \text{ProjBias}_l(p) = \mathbf{s}_l(p) \cdot \mathbf{g}_l
  \label{eq:proj}
\end{equation}
This is the scalar projection of the (un-normalised) profession vector onto the
unit gender direction. It preserves vector magnitude, so professions whose
representations lie further from the origin receive proportionally higher
scores.

\paragraph{Cosine difference (CosDiff).}
\begin{equation}
  \text{CosDiff}_l(p) = \cos\!\left(\mathbf{s}_l(p),\, \mathbf{m}_l\right)
                       - \cos\!\left(\mathbf{s}_l(p),\, \mathbf{f}_l\right)
  \label{eq:cosdiff}
\end{equation}
This is a purely angular measure, invariant to vector magnitude. A positive
value indicates that $p$'s representation is closer (in angle) to the male
centroid than to the female centroid.

Both estimators are computed for every layer $l \in \{0, 1, \ldots, 12\}$ and
for every profession in the set of 30 Tigrigna job titles. Their
layer-wise Spearman rank correlation tests whether the two metrics agree on the
ordering of professions by bias magnitude (\textbf{RQ3}).

\subsection{Representation Uniqueness Analysis}

To assess whether contextual embeddings genuinely encode sentence-level context
or collapse to near-static representations, we compute the \emph{mean pairwise
cosine similarity} across all sentence vectors for each profession word. For a
profession appearing in $N$ sentences, we extract $N$ layer-$l$ vectors and
compute:
\begin{equation}
  \overline{\cos}_l(p) = \frac{2}{N(N-1)}
    \sum_{i < j} \cos\!\left(\mathbf{v}^{(i)}_l,\, \mathbf{v}^{(j)}_l\right)
\end{equation}
A value near 1.0 indicates that the model produces nearly identical
representations regardless of sentence context (\emph{static} regime); a lower
value indicates context-sensitivity.

To disentangle corpus template effects from model-intrinsic behaviour, we also
ran the analysis on five hand-crafted sentences for \textit{ሓኪም} (doctor)
spanning maximally diverse syntactic roles (subject/agent, object/recipient,
negated goal, embedded relative clause, hypothetical conditional).

\subsection{Static Embedding Analysis}

As a complementary probe, we extract each profession word's \emph{static
embedding} directly from the model's token embedding weight matrix
$\mathbf{W} \in \mathbb{R}^{V \times 768}$, with no forward pass. This gives
the model's a-priori representation before any transformer attention is
applied:
\begin{equation}
  \mathbf{e}(p) = \frac{1}{k} \sum_{i=1}^{k} \mathbf{W}[t_i]
\end{equation}
We build an analogous static gender direction from the anchor words' static
embeddings and project each profession's static vector onto it. Comparing
the static bias score with the contextual bias score reveals whether gender
association is encoded in the model weights themselves or emerges from the
transformer layers.

\subsection{Batch Stability Validation}

To verify that bias estimates are robust to sentence sampling, we partition
each profession's 499 corpus sentences into 24 non-overlapping batches of
20 sentences each (19 sentences are unused as a leftover). We compute the mean
projection score independently for each batch and report the grand mean,
standard deviation, minimum, and maximum across the 24 batches. A low standard
deviation confirms that the bias estimate is sampling-independent.


% ============================================================
\section{Results}
% ============================================================

\subsection{Representation Uniqueness}

Table~\ref{tab:mean_cos} reports the mean pairwise cosine similarity
$\overline{\cos}$ (averaged across all 13 layers) for every profession word,
computed from 50 sentences per profession.

\begin{table}[h]
\centering
\small
\begin{tabular}{lcc}
\toprule
\textbf{Profession} & \textbf{$\overline{\cos}$ (all layers)} & \textbf{Subword tokens} \\
\midrule
ኣካውንታንት  (accountant)     & 0.9809 & 5 \\
ሰራሕተኛ ጽሬት (cleaner)       & 0.9812 & 6 \\
ኣርኪቴክት   (architect)      & 0.9803 & 5 \\
ኤለክትሪሻን  (electrician)    & 0.9782 & 6 \\
ዳይረክተር   (director)       & 0.9772 & 5 \\
ሓላዊ ጸጥታ  (security guard)  & 0.9775 & 5 \\
\midrule
\multicolumn{3}{c}{\ldots} \\
\midrule
ጠበቓ       (lawyer)        & 0.9196 & 3 \\
\midrule
\textbf{Mean (all 30)}     & \textbf{0.969} & --- \\
\textbf{Min}               & \textbf{0.920} (ጠበቓ) & --- \\
\textbf{Max}               & \textbf{0.981} (ሰራሕተኛ ጽሬት) & --- \\
\bottomrule
\end{tabular}
\caption{Mean pairwise cosine similarity across sentence vectors (50 sentences,
all 13 layers). All values exceed 0.91, indicating near-static contextual
representations. ጠበቓ (lawyer) has the lowest value and is the most
context-sensitive profession in the set.}
\label{tab:mean_cos}
\end{table}

All 30 professions exceed $\overline{\cos} = 0.91$, indicating that the model
produces near-identical representations for every profession word regardless of
sentence context. For \textit{ሓኪም} (doctor) specifically, we observe
$\overline{\cos} = 0.967$ across all 499 corpus sentences
($124{,}251$ pairwise comparisons at layer~12 alone; 100\% of pairs have
cosine~$> 0.99$ at that layer).

To test whether this is a corpus-template artifact, we ran the same analysis
on five hand-crafted sentences for \textit{ሓኪም} spanning maximally diverse
syntactic and semantic contexts. The result was $\overline{\cos} = 0.947$
— a reduction of only 0.020 compared to the template corpus — confirming that
near-static representations are an intrinsic property of the model's handling
of Tigrigna tokens, not an artifact of corpus construction.

\subsection{Layer-wise Projection Curve}

Figure~\ref{fig:layer_curve} shows the mean projection score across all 30
professions at each transformer layer.

\begin{figure}[h]
  \centering
  \includegraphics[width=\linewidth]{output/ti/xlm-roberta-base/figs/ti_projection_curve.png}
  \caption{Layer-wise mean projection score (ProjBias) averaged across 30
  Tigrigna professions on XLM-RoBERTa-base. Positive values indicate
  male-leaning bias. The curve is near-neutral at layer~0 (+0.20), dips
  slightly negative at layers~2--4 (trough: $-0.62$ at layer~2), then rises
  sharply to a peak of $+2.99$ at layer~9, before collapsing to near-zero
  ($+0.04$) at the final layer~12.}
  \label{fig:layer_curve}
\end{figure}

The curve reveals a consistent three-phase structure:
\begin{enumerate}
  \item \textbf{Embedding phase (layers 0--4):} Representations are near-neutral
    or slightly female-leaning. The model is primarily encoding surface-level
    token features. Mean projection ranges from $+0.20$ (layer~0) to
    $-0.62$ (layer~2).
  \item \textbf{Semantic accumulation phase (layers 5--9):} A monotonic rise in
    male bias, from $+0.52$ (layer~5) to a peak of $+\mathbf{2.99}$ at
    \textbf{layer~9}. This is where the transformer's attention heads build
    semantic representations and gender associations accumulate.
  \item \textbf{Output collapse phase (layers 10--12):} A sharp reduction,
    with the final layer~12 reaching $+0.04$ — effectively neutral. This is
    consistent with the representation uniqueness finding that layer~12 produces
    fully anisotropic, near-identical vectors ($\overline{\cos}_{12} = 0.997$
    for \textit{ሓኪም}).
\end{enumerate}

\subsection{Metric Agreement — Spearman Correlation}

Table~\ref{tab:spearman} reports the Spearman rank correlation $\rho$ between
ProjBias and CosDiff across all 30 professions at each layer.

\begin{table}[h]
\centering
\small
\begin{tabular}{ccc}
\toprule
\textbf{Layer} & \textbf{Spearman $\rho$} & \textbf{$p$-value} \\
\midrule
0  & 0.9947 & $< 10^{-28}$ \\
1  & 0.9929 & $< 10^{-27}$ \\
2  & 0.9933 & $< 10^{-27}$ \\
3  & 0.9938 & $< 10^{-27}$ \\
4  & 0.9929 & $< 10^{-27}$ \\
5  & 0.9978 & $< 10^{-33}$ \\
6  & 0.9969 & $< 10^{-31}$ \\
7  & 0.9924 & $< 10^{-26}$ \\
8  & 0.9947 & $< 10^{-28}$ \\
9  & 0.9956 & $< 10^{-29}$ \\
10 & 0.9987 & $< 10^{-36}$ \\
11 & 1.0000 & $0$ \\
12 & 1.0000 & $0$ \\
\bottomrule
\end{tabular}
\caption{Spearman rank correlation between ProjBias and CosDiff across 30
professions at each layer ($n=30$, all $p < 10^{-26}$). Both metrics agree
perfectly at layers~11--12 ($\rho = 1.00$) and near-perfectly at all
other layers ($\rho \geq 0.992$).}
\label{tab:spearman}
\end{table}

The two estimators agree to an extraordinary degree: the minimum Spearman
$\rho$ across all 13 layers is $0.992$ (layer~7), rising to a perfect
$\rho = 1.00$ at layers~11 and~12. All correlations are significant at
$p < 10^{-26}$. This near-perfect agreement validates that both metrics
measure the same underlying bias structure, and that the choice of estimator
does not affect the rank ordering of professions by gender bias.

\subsection{Profession-level Bias Analysis}

Table~\ref{tab:prof_bias} summarises the mean ProjBias (averaged across all 13
layers) for key professions, alongside their static embedding bias scores.

\begin{table}[h]
\centering
\small
\begin{tabular}{lcccl}
\toprule
\textbf{Profession} & \textbf{Static proj} & \textbf{Contextual mean} & \textbf{Amplification} & \textbf{Mechanism} \\
\midrule
መምህር  (teacher)  & $+0.499$ & $+2.084$ & $+1.585$ & Embedding-level \\
ፖሊስ   (police)   & $+0.428$ & ---      & ---       & Embedding-level \\
ሓረስታይ (farmer)  & $+0.199$ & $+1.123$ & $+0.924$ & Attention-level \\
ሓኪም   (doctor)  & $+0.012$ & $+0.640$ & $+0.628$ & Attention-level \\
ጠበቓ   (lawyer)  & $-0.259$ & $-0.103$ & ---       & Consistent neutral \\
ነርስ    (nurse)   & $-0.352$ & $+2.748$* & $+3.100$ & Attention-level (reversal) \\
ነዳቓይ  (miner)   & $-0.489$ & ---       & ---       & Statically female \\
ሳይንቲስት (scientist) & $-0.508$ & ---     & ---       & Statically female \\
\bottomrule
\end{tabular}
\caption{Static vs.\ contextual mean projection bias for selected professions.
Amplification = contextual $-$ static. *Nurse contextual value is the mean
ProjBias across layers 0--12 from the bulk run; the peak at layer~9 is
$+4.29$.}
\label{tab:prof_bias}
\end{table}

\subsection{Batch Stability Validation}

Table~\ref{tab:batch_stability} reports the grand mean, standard deviation,
and range of the per-batch mean projection score across 24 batches of
20 sentences each, for four professions.

\begin{table}[h]
\centering
\small
\begin{tabular}{lccccl}
\toprule
\textbf{Profession} & \textbf{Grand mean} & \textbf{Std dev} & \textbf{Min} & \textbf{Max} & \textbf{Batches (M/F/N)} \\
\midrule
መምህር  (teacher)  & $+2.084$ & $0.045$ & $+2.009$ & $+2.181$ & 24 / 0 / 0 \\
ሓረስታይ (farmer)  & $+1.123$ & $0.067$ & $+0.998$ & $+1.248$ & 24 / 0 / 0 \\
ሓኪም   (doctor)  & $+0.640$ & $0.092$ & $+0.529$ & $+0.850$ & 24 / 0 / 0 \\
ጠበቓ   (lawyer)  & $-0.103$ & $0.089$ & $-0.230$ & $+0.091$ & ~0 / 0 / 24 \\
\bottomrule
\end{tabular}
\caption{Batch stability analysis across 24 batches of 20 sentences each.
M/F/N counts indicate how many batches received a MALE-LEANING / FEMALE-LEANING
/ ROUGHLY NEUTRAL verdict. All 24 batches for teacher, farmer, and doctor
are MALE-LEANING regardless of the local male/female sentence ratio.}
\label{tab:batch_stability}
\end{table}

For \textit{መምህር} (teacher), the standard deviation is only $0.045$ — the
lowest of any profession tested — meaning the per-batch estimates vary by less
than 2\% of the grand mean. Even batches drawn from predominantly female-context
sentences return scores above $+2.0$. The floor for teacher ($+2.009$, batch~20)
is higher than the \emph{ceiling} for doctor ($+0.850$, batch~02), demonstrating
that these two professions occupy entirely non-overlapping ranges of male bias.

For \textit{ጠበቓ} (lawyer), no batch crosses the $\pm 0.3$ threshold in
either direction (range: $-0.230$ to $+0.091$), confirming consistent
neutrality independent of sampling.


% ============================================================
\section{Observations and Discussion}
% ============================================================

\subsection{Near-static Representations in a Low-Resource Language}

Our first and most foundational finding is that XLM-RoBERTa-base produces
effectively static representations for all 30 Tigrigna profession tokens.
The mean pairwise cosine similarity across sentence vectors exceeds $0.91$ for
every profession (Table~\ref{tab:mean_cos}), indicating that the model assigns
nearly identical hidden-state vectors to a given token regardless of the
surrounding discourse context.

This is not primarily a corpus artifact. The custom-sentence experiment
for \textit{ሓኪም}, using five structurally and semantically maximally diverse
sentences (active agent, object-recipient, negated goal, embedded relative
clause, hypothetical conditional), produced $\overline{\cos} = 0.947$ compared
to $0.967$ from the template corpus — a reduction of only $0.020$. The near-static
behaviour is therefore intrinsic to the model's treatment of Tigrigna tokens,
not an artifact of repetitive sentence structure.

We attribute this to the low-resource nature of Tigrigna within XLM-RoBERTa's
pretraining corpus. When a language is sparsely represented in pretraining data,
the model's attention heads have insufficient statistical signal to learn
context-sensitive representations for that language's tokens. The result is
contextual collapse: the transformer degenerates to a near-lookup-table
behaviour, producing representations that are driven primarily by the token
identity rather than by surrounding context.

\subsection{Two Distinct Mechanisms of Gender Bias}

Comparing static embedding bias scores with contextual bias scores reveals that
gender bias in Tigrigna profession words arises through two qualitatively
distinct mechanisms:

\paragraph{Type 1: Embedding-level bias.}
For \textit{መምህር} (teacher, static proj $= +0.499$, contextual $= +2.084$),
the male association is already encoded in the raw token embedding — before
any attention is applied. The transformer layers then amplify this signal
further, producing a contextual bias more than four times larger than the
static score. The bias originates in the pretraining data distribution and is
present even at the input representation level.

\paragraph{Type 2: Attention-level bias.}
For \textit{ሓኪም} (doctor, static proj $= +0.012$) and \textit{ሓረስታይ}
(farmer, static $= +0.199$), the raw token embeddings are essentially gender-neutral.
Yet the transformer layers impose a strong male bias (+0.640 and +1.123,
respectively) during forward propagation. The attention mechanism is introducing
gender associations that are absent from the input representation — presumably
by activating co-occurrence patterns learned from the broader multilingual
pretraining data in which the model has richer coverage.

\paragraph{The nurse reversal.}
The most dramatic instance of Type~2 bias is \textit{ነርስ} (nurse), which has a
statically \emph{female}-leaning embedding (static proj $= -0.352$) but
emerges from the transformer with a strongly \emph{male}-leaning contextual
representation, peaking at $+4.29$ at layer~9 — the highest single-profession
value in our entire dataset. The transformer not only overrides but completely
reverses the static signal. This suggests that the multilingual pretraining
signal for nurse in other languages (where the model has richer data) is
strongly male-associated, overriding the Tigrigna token's own embedding.

\subsection{Layer-wise Bias Structure}

The layer-wise projection curve (Figure~\ref{fig:layer_curve}) exhibits a
consistent three-phase pattern. Early layers (0--4) are near-neutral or
slightly female-leaning, reflecting morphological token processing. Middle-to-deep
layers (5--9) show a monotonic accumulation of male bias peaking at layer~9
($+2.99$ mean), where semantic representations are most fully formed.
The final layer (12) collapses to near-zero ($+0.04$), consistent with the
observation that layer~12 produces the most anisotropic, context-collapsed
representations ($\overline{\cos}_{12} \approx 0.997$ for \textit{ሓኪም}).

Layer~9 is therefore the critical layer for gender bias measurement in
XLM-RoBERTa on Tigrigna. Individual profession scores at this layer span
a wide range: from $+5.90$ for \textit{መምህር} (teacher) to $-1.27$ for
\textit{ነርስ} at layer~0, with the profession ordering preserved near-perfectly
across layers (Spearman $\rho \geq 0.992$, Table~\ref{tab:spearman}).

\subsection{Lawyer as a Contrast Case}

\textit{ጠበቓ} (lawyer) is the single profession that remains consistently
neutral across all analyses. It has the lowest mean pairwise cosine similarity
of all 30 professions ($\overline{\cos} = 0.920$), indicating the most
contextually sensitive representations. It is the most statically female-adjacent
among the neutral-labelled words (static proj $= -0.259$), yet the transformer
does not amplify this signal. All 24 sampling batches return a ROUGHLY NEUTRAL
verdict (grand mean $= -0.103$, range $[-0.230, +0.091]$).

This combination — highest context-sensitivity, lowest contextual bias — is
consistent with the hypothesis that contextual diversity suppresses bias
amplification. When the model's attention has more to work with, it does not
default to gender stereotyping. The inverse relationship between
$\overline{\cos}$ and bias magnitude, observed across the profession set,
suggests that representation diversity and gender neutrality are positively
correlated in this setting.

\subsection{Implications for Low-Resource Language Bias Research}

Our results carry several implications for the study of gender bias in
multilingual language models, particularly for low-resource languages:

\begin{enumerate}
  \item \textbf{Contextual bias measurement may be unreliable for low-resource
    languages.} When representations are near-static ($\overline{\cos} > 0.95$),
    the averaging step in contextual embedding extraction adds negligible
    information. The bias score is effectively being read from a single vector,
    making it equivalent to probing the static embedding — but without the
    explicitness of doing so directly.

  \item \textbf{Batch stability validates the measurement.} The consistently
    low standard deviation across 24 sampling batches (as low as $0.045$ for
    teacher) demonstrates that bias estimates derived from 20 sentences are
    already fully converged for near-static representations. This provides a
    principled basis for using small profession context windows in low-resource
    settings.

  \item \textbf{The two-mechanism framework enables targeted debiasing.}
    Type~1 bias (embedding-level) may be addressable through targeted
    embedding re-initialisation or gender-neutral retrofitting of the static
    weight matrix. Type~2 bias (attention-level) requires intervention in the
    transformer computation itself — for example, through projection-based
    nullspace methods applied to the attention outputs at the peak bias layer
    (layer~9 in our case).

  \item \textbf{Perfect metric agreement simplifies reporting.} The near-perfect
    Spearman correlation ($\rho \geq 0.992$) between ProjBias and CosDiff at
    every layer means that both metrics are interchangeable for ranking
    professions by bias magnitude in this setting. Researchers can safely
    report either metric without loss of comparative validity.
\end{enumerate}
